\section{Studi Kasus: Optimasi Kode Bercabang \textit{If-Else} dan Gerbang Logika}

Penerapan aljabar logika tidak berhenti pada pembuktian teoritis. Dalam praktik, aljabar logika digunakan untuk menyederhanakan percabangan program (\textit{software}) dan mengoptimalkan rancangan gerbang digital (\textit{hardware}) \cite{rosen2019}.

\subsection{Penyederhanaan Percabangan \textit{If-Else}}

Ekspresi logika yang terlalu kompleks dalam kode dapat menyulitkan pemeliharaan dan meningkatkan risiko kesalahan.

\textbf{Skenario Kode Asli Mentah:}
\begin{lstlisting}[style=pythonstyle]
def cek_validitas_pinjaman(a_kerja, b_jaminan, c_bebas_utang):
  
  # Validasi kualifikasi nasabah (versi ruwet):
  if not (not a_kerja and b_jaminan):
      # Cek tambahan di dalam:
      if a_kerja and (b_jaminan or not c_bebas_utang):
          return True
  return False
\end{lstlisting}

Representasi logikanya:
\[
\neg(\neg A \wedge B) \wedge (A \wedge (B \vee \neg C))
\]

\textbf{Langkah penyederhanaan:}
\begin{align*}
    1.\ & \neg(\neg A \wedge B) \wedge A \wedge (B \vee \neg C) & \text{\parbox[t]{5.8cm}{Bentuk awal}} \\
    2.\ & (A \vee \neg B) \wedge A \wedge (B \vee \neg C) & \text{\parbox[t]{5.8cm}{De Morgan}} \\
    3.\ & A \wedge (A \vee \neg B) \wedge (B \vee \neg C) & \text{\parbox[t]{5.8cm}{Komutatif}} \\
    4.\ & \mathbf{A} \wedge (B \vee \neg C) & \text{\parbox[t]{5.8cm}{Absorpsi}}
\end{align*}

\textbf{Transformasi Kode Bersih Pasca-Penyulingan (Refactoring):}
\begin{lstlisting}[style=pythonstyle]
def cek_validitas_pinjaman(a_kerja, b_jaminan, c_bebas_utang):
  
  # Filter Ekuivalensi Ceria nan Minimalis:
  return a_kerja and (b_jaminan or not c_bebas_utang)
\end{lstlisting}
Hasil refactoring menghasilkan logika yang lebih ringkas, lebih mudah dibaca, dan lebih mudah diuji.

\subsection{Aplikasi pada Rangkaian Gerbang Logika Digital}

Logika proposisional juga menjadi dasar optimasi rangkaian digital. Misalkan rancangan sinyal mengikuti ekspresi:
$$ Q_{output} = (A \land B) \lor (A \land \neg B) $$

\textbf{Penyederhanaan:}
\begin{align*}
    (A \wedge B) \vee (A \wedge \neg B) &\equiv A \wedge (B \vee \neg B) \quad \text{(Distributif)} \\
    &\equiv A \wedge \mathbf{T} \quad \text{(Hukum Negasi)} \\
    &\equiv \mathbf{A} \quad \text{(Identitas)}
\end{align*}

Dengan demikian, $Q_{output} \equiv A$. Artinya, rangkaian dapat disederhanakan secara signifikan karena keluaran cukup mengikuti sinyal $A$ \cite{wiki_proplogic}.
